\section{Communications (802.11p, VLC, LTE/5G)}
\label{sec:sota_communications}

V2X communications for platooning must sustain periodic and low-latency exchanges, including beaconing and control messages, in highly dynamic and often congested environments. Unlike infotainment applications or simple information dissemination, here the effect of a loss or a delay immediately propagates to the controller dynamics. For this reason, metrics such as PDR/FER, latency, and jitter must be interpreted in relation to the role played by the technology within the control chain, rather than as isolated indicators \cite{segata2015communication,dressler2019cooperative,segata2023multi-technology}.

IEEE~802.11p is the historical reference for dedicated short- and medium-range vehicular communications. The standard was conceived as an adaptation of the IEEE~802.11 family to vehicular conditions, with 10~MHz bandwidth, transmission over multiple subcarriers in parallel, and modifications designed to cope with the mobile channel and support direct communication without traditional association to access points \cite{ieee80211p-2010,segata2016safe}. In the platooning domain, 802.11p has been extensively studied both in simulation and experimentally, showing good performance in scenarios that are not excessively congested but also critical weaknesses in the presence of high vehicle density, competing traffic, and random medium access \cite{segata2015communication}. In particular, the contention-based nature of the MAC can produce burst losses and irregular channel access times, with a significant impact on the most demanding CACC schemes.

Visible Light Communication (VLC) has been proposed as a complementary technology for short-range intra-platoon communications, exploiting directionality and line of sight between vehicle headlights and taillights. The literature shows that integrating VLC with 802.11p can reduce radio load and increase resilience in specific scenarios, especially when the platoon geometry favors stable links between adjacent vehicles \cite{ishihara2015improving,segata2016platooning,schettler2019deeply}. At the same time, VLC has structural limitations: dependence on optical visibility, sensitivity to occlusions and misalignment, and limited useful range. The result is a reliability profile that complements, rather than replaces, radio technologies \cite{hardes2019communication}.

Cellular technologies, namely LTE/C-V2X and, more recently, 5G, instead provide broader coverage, infrastructure support, and the possibility of integration with advanced network services. In the multi-technology work of Segata \emph{et al.}, the cellular interface considered is LTE C-V2X, Mode~3 under network coverage, and is therefore subject to handover phenomena and infrastructure dependence \cite{segata2023multi-technology,3gpp-21914-1400}. More generally, cellular networks introduce important opportunities, such as service continuity, coordination, and managed resources, but also specific operational weaknesses, such as loss of line of sight, variability in link quality, and transients due to handovers \cite{bazzi2020wireless}.

In summary, 802.11p, VLC, and LTE/5G exhibit different and only partially overlapping performance and failure profiles: 802.11p mainly suffers from congestion, interference, and contention; VLC offers high directionality but comes with LOS and range constraints; the cellular network extends coverage but introduces infrastructure dependencies and handover transients.
